\documentclass[12pt,a4paper]{article}

\usepackage[utf8]{inputenc}
\usepackage[T1]{fontenc}
\usepackage[polish]{babel}
\usepackage{amsmath}
\usepackage{amssymb}
\usepackage{geometry}
\usepackage{array}
\usepackage{listings}
\usepackage{needspace}

\geometry{margin=2.5cm}

% Naglowki sekcji nie powinny zostawac same na koncu strony.
\newcommand{\sekcja}[1]{%
  \Needspace{8\baselineskip}%
  \section*{#1}%
  \nopagebreak[4]%
}

% Blok naglowek + krotki wstep trzyma razem na jednej stronie.
\newcommand{\sekcjazwstep}[2]{%
  \Needspace{10\baselineskip}%
  \section*{#1}%
  \nopagebreak[4]%
  #2%
}

\widowpenalty=10000
\clubpenalty=10000

\lstset{
    basicstyle=\ttfamily\small,
    breaklines=true,
    frame=single,
    columns=fullflexible,
    keepspaces=true,
    showstringspaces=false
}

\title{Podstawy Symulacji Komputerowej\\Lista 2 -- zadanie 5}
\author{}
\date{}

\begin{document}

\maketitle

\sekcjazwstep{Treść zadania}{%
Metodą zdarzeń dyskretnych rozwiązać następujące zadanie:
}

\begin{quote}
Zakładamy, że szkody klientów zgłaszane są według procesu Poissona
o intensywności 10 na dzień. Kwota szkody jest zmienną losową o rozkładzie
wykładniczym ze średnią \$1000. Firma ubezpieczeniowa otrzymuje wpłaty
w sposób ciągły ze stałą intensywnością \$11000 na dzień. Przyjmując kapitał
początkowy \$25000, oszacować metodą symulacji prawdopodobieństwo, że kapitał
pozostanie dodatni przez pierwsze 365 dni.
\end{quote}

\sekcjazwstep{Model probabilistyczny}{%
Niech $U(t)$ oznacza kapitał firmy w chwili $t$, gdzie $t\in[0,365]$.
}
Przyjmujemy:
\[
U(0)=25000.
\]

Wpłaty napływają ciągle ze stałą intensywnością:
\[
c=11000 \quad \text{[\$ na dzień]}.
\]

Między kolejnymi szkodami kapitał zmienia się liniowo:
\[
\frac{\mathrm{d}U}{\mathrm{d}t}=c,
\]
o ile nie wystąpi zdarzenie szkody.

Zgłoszenia szkód tworzą proces Poissona o intensywności:
\[
\lambda=10 \quad \text{[zdarzeń na dzień]}.
\]

Oznacza to, że czasy między kolejnymi szkodami są niezależne i mają rozkład
wykładniczy:
\[
T_i \sim \mathrm{Exp}(\lambda),
\]
przy czym średni czas między szkodami wynosi $1/\lambda=0{,}1$ dnia.

Kwota pojedynczej szkody $X_i$ jest zmienną losową o rozkładzie wykładniczym
ze średnią \$1000:
\[
E(X_i)=1000,
\]
czyli parametr skali wynosi $1000$, a gęstość:
\[
f(x)=\frac{1}{1000}e^{-x/1000},\qquad x>0.
\]

W chwili wystąpienia $i$-tej szkody kapitał ulega natychmiastowej redukcji:
\[
U(t_i^+)=U(t_i^-)-X_i.
\]

Szukane jest prawdopodobieństwo:
\[
p=P\big(U(t)>0 \text{ dla wszystkich } t\in[0,365]\big).
\]

\sekcjazwstep{Metoda zdarzeń dyskretnych}{%
W metodzie zdarzeń dyskretnych symulujemy tylko chwile, w których stan systemu
ulega zmianie. W tym modelu są to wyłącznie chwile zgłoszenia szkód.

Algorytm jednej replikacji:
}

\Needspace{16\baselineskip}
\begin{enumerate}
    \item Ustaw $t=0$ oraz $U=25000$.
    \item Wygeneruj odstęp do kolejnej szkody:
    \[
    \Delta \sim \mathrm{Exp}(\lambda).
    \]
    \item Jeżeli $t+\Delta>365$, dolicz wpłaty do końca roku i zakończ replikację.
    \item W przeciwnym razie:
    \begin{itemize}
        \item zwiększ kapitał o wpłaty w okresie $(t,t+\Delta)$:
        \[
        U \leftarrow U + c\Delta;
        \]
        \item ustaw $t\leftarrow t+\Delta$;
        \item wygeneruj kwotę szkody $X\sim \mathrm{Exp}(1000)$;
        \item zmniejsz kapitał: $U\leftarrow U-X$;
        \item jeżeli $U\le 0$, zakończ replikację wynikiem negatywnym.
    \end{itemize}
    \item Powtarzaj krok 2.
\end{enumerate}

Ponieważ między szkodami kapitał rośnie liniowo, minimum na każdym takim
odcinku występuje bezpośrednio po obsłużeniu szkody. Wystarczy zatem sprawdzać
warunek dodatniego kapitału po każdym zdarzeniu szkody.

\sekcjazwstep{Implementacja}{%
Symulację wykonano w języku Python. Poniżej przedstawiono główną funkcję
jednej replikacji:
}

\Needspace{22\baselineskip}
\begin{lstlisting}[language=Python]
def symuluj_jedna_sciezke():
    czas = 0.0
    kapital = 25000.0

    while True:
        odstep = random.expovariate(10.0)
        czas_szkody = czas + odstep

        if czas_szkody > 365.0:
            kapital += 11000.0 * (365.0 - czas)
            return kapital > 0.0

        kapital += 11000.0 * (czas_szkody - czas)
        czas = czas_szkody

        kwota_szkody = random.expovariate(1.0 / 1000.0)
        kapital -= kwota_szkody

        if kapital <= 0.0:
            return False
\end{lstlisting}

Prawdopodobieństwo $p$ estymujemy metodą Monte Carlo:
\[
\hat{p}=\frac{1}{N}\sum_{j=1}^{N} I_j,
\]
gdzie $I_j=1$, jeżeli w $j$-tej replikacji kapitał pozostał dodatni przez cały
okres $[0,365]$, a w przeciwnym razie $I_j=0$.

W obliczeniach przyjęto:
\[
N=100000.
\]

\sekcjazwstep{Wyniki symulacji}{%
Po wykonaniu $N=100000$ niezależnych replikacji otrzymano:
}

\Needspace{10\baselineskip}
\[
\begin{array}{l|r}
\text{Wielkość} & \text{Wartość}\\
\hline
\text{Liczba replikacji} & 100000\\
\text{Liczba sukcesów} & 90637\\
\text{Liczba niepowodzeń} & 9363\\
\text{Estymator } \hat{p} & 0{,}9064\\
\text{95\% przedział ufności} & (0{,}9045;\ 0{,}9082)
\end{array}
\]

Oznacza to, że w około $90{,}64\%$ symulowanych scenariuszy kapitał firmy
pozostał dodatni przez cały rok.

\sekcjazwstep{Interpretacja wyniku}{%
Średnia liczba szkód w ciągu 365 dni wynosi około:
}
\[
\lambda \cdot 365 = 10 \cdot 365 = 3650.
\]

Średni łączny koszt szkód:
\[
3650 \cdot 1000 = 3\,650\,000.
\]

Średni napływ wpłat w ciągu roku:
\[
11000 \cdot 365 = 4\,015\,000.
\]

Teoretyczna średnia zmiana kapitału (bez uwzględnienia zmienności) jest dodatnia,
jednak duża losowość kwot pojedynczych szkód powoduje, że w części scenariuszy
następuje bankructwo.

\sekcjazwstep{Odpowiedź końcowa}{%
Metodą zdarzeń dyskretnych oszacowano, że prawdopodobieństwo utrzymania dodatniego
kapitału przez pierwsze 365 dni wynosi:
}

\Needspace{12\baselineskip}
\[
\boxed{\hat{p}\approx 0{,}9064}
\]

czyli około:

\[
\boxed{90{,}6\%}.
\]

Przy $95\%$ poziomie ufności można przyjąć przedział:

\[
\boxed{(0{,}9045;\ 0{,}9082)}.
\]

\end{document}
